\chapter{Conclusion and Future Work}
\label{ch:conclusion}

Draw your conclusions from the results and summarize your contributions.
Point out aspects that need to be investigated further.

The conclusion can be structured inversely to the introduction:
Summarize how the \emph{evaluation} backs the \emph{solution}, which solves the \emph{problem}.
Describe how your contributions improve the \emph{situation} and what other (potentially newly discovered) problems have to be solved in the future.

Be concise: the conclusion normally fits on a single page and is rarely longer than two pages.


\section{Summary}
\label{sec:conclusion_summary}

\section{Future work}
\label{sec:conclusion_future_work}

\subsection{Model accuracy}
\label{subsec:future_accuracy}
% oversampled nonlinearity (L = 4-8), more branches, structured WH initialization,
% knee-resolving spline with linear extrapolation, per-level loss weighting,
% lower measurement levels for the hard setting, Python refinement on near-linear systems

\subsection{Hardware}
\label{subsec:future_hardware}
% calibration of the Grätli chain (PRE/POST GAIN, bypass model),
% code optimization (state in locals, fused passes, CMSIS-DSP, 480 MHz), stereo

\subsection{Evaluation}
\label{subsec:future_evaluation}
% ABX listening test with several listeners, band-limited test signals